% Generated draft: requires mathematical and source review.
\documentclass[11pt]{article}
\usepackage[margin=25mm]{geometry}
\usepackage{fontspec}
\setmainfont{Latin Modern Roman}
\usepackage{amsmath,amssymb,mathtools,unicode-math}
\setmathfont{Latin Modern Math}
\usepackage{enumitem,needspace,xurl,hyperref}
\setlength{\parindent}{0pt}
\setlength{\parskip}{5pt}
\setlength{\emergencystretch}{3em}
\newcommand{\N}{\mathbb{N}}
\newcommand{\Z}{\mathbb{Z}}
\newcommand{\Q}{\mathbb{Q}}
\newcommand{\R}{\mathbb{R}}
\newcommand{\C}{\mathbb{C}}
\newcommand{\F}{\mathbb{F}}
\newcommand{\A}{\mathbb{A}}
\newcommand{\PP}{\mathbb{P}}
\newcommand{\E}{\mathbb{E}}
\newcommand{\eps}{\varepsilon}
\newcommand{\dd}{\,\mathrm{d}}
\begin{document}
\section*{UnsolvedMath ID: 30002754}
\textbf{Source problem code:} catalogue entry
\subsection{Definitions and mathematical statement}
Write \(H^1(\Omega)\) for the first-order Sobolev space of square-integrable functions on \(\Omega\) whose first weak derivatives are square integrable. The supplied record does not specify the ambient dimension, scalar field, or standing hypotheses on the domain.

The supplied literature assessment identifies the relevant notion as an approximative trace, but the record does not define that notion, its boundary measure, its target space, or its convergence requirements. To state the questions without selecting these missing data, let \(B_\Omega\) denote the intended boundary-value space and write
\[
\mathcal T_\Omega(u)=\{\varphi\in B_\Omega:\varphi\text{ is a trace of }u\text{ in the intended source sense}\}.
\]
This is notation for the unspecified source trace relation, not a replacement definition of approximative trace. Existence means \(\mathcal T_\Omega(u)\ne\varnothing\); having at most one trace means
\[
\forall\varphi,\psi\in\mathcal T_\Omega(u),\qquad \varphi=\psi
\quad\text{in }B_\Omega.
\]
Let \(V_\Omega=\widetilde H^1(\Omega)\subseteq H^1(\Omega)\) denote the smaller space named in the record. Its defining closure, approximation class, and topology are not supplied.

The first question is whether, for every domain \(\Omega\) with continuous boundary in the intended ambient class,
\[
\forall u\in H^1(\Omega),\qquad
\forall\varphi,\psi\in\mathcal T_\Omega(u),\quad \varphi=\psi.
\tag{1}
\]
Here ``continuous boundary'' is retained from the source; its precise local geometric definition and any additional domain hypotheses are unavailable in the supplied record.

The explicitly stated particular question asks whether the following implication holds for every eligible domain \(\Omega\):
\[
\left[\forall v\in\widetilde H^1(\Omega),\quad
\mathcal T_\Omega(v)\ne\varnothing\right]
\quad\Longrightarrow\quad
\left[\forall u\in H^1(\Omega),\quad
\forall\varphi,\psi\in\mathcal T_\Omega(u),\quad\varphi=\psi\right].
\tag{2}
\]
The record also asks whether existence of a trace implies uniqueness, but does not specify whether this phrase concerns an individual function or a universal existence property on a function space. Equation (2) preserves its expressly quantified instance. It is unclear whether the domain class for (2) retains the continuous-boundary hypothesis from (1) or is broader. These omissions prevent a fully specified analytic formulation.

The supplied dated assessment labels the problem partially solved: it attributes an affirmative answer to (1) to Sauter and says that the implication in (2) is explicitly left unresolved. These are supplied literature claims, not independently verified conclusions.
\subsection{Short English statement}
The questions concern uniqueness of approximative Sobolev traces on domains with continuous boundary, and whether existence of traces for every function in \(\widetilde H^1(\Omega)\) forces uniqueness for every function in \(H^1(\Omega)\). The supplied assessment reports the first question answered and the implication unresolved. Precise trace definitions, domain hypotheses, and the scope of the second question require source review.
\subsection{Sources}
\begin{itemize}
\item Supplied local UnsolvedMath registry record, permanent ID \(30002754\), title ``Existence and Uniqueness of Sobolev Traces on Continuous Domains,'' exact source problem code \texttt{catalogue entry}, in the Oberwolfach Reports Open Problems collection. The record gives proposed year 2015 and names no proposer. Its labels \texttt{partially\_solved} and \texttt{corrected\_verified}, including the claim that notation, typography, or omitted definitions were repaired while preserving mathematical content, are supplied registry assertions.
\item \emph{Spectral Theory and Weyl Functions} (2015), Oberwolfach Reports, pp.\ 5--86, DOI 10.4171/owr/2015/1. The record describes these reports as documenting mathematical workshops and their open-problem discussions. Supplied source report and literature-source link: \url{https://doi.org/10.4171/owr/2015/1}.
\item Supplied literature-source label: ``Sauter, uniqueness of the approximative trace.'' \url{https://arxiv.org/abs/1607.07109}. The supplied assessment attributes continuous-boundary uniqueness to this work and says it explicitly leaves the smaller-space existence implication unresolved.
\item Literature checked at: \texttt{2026-08-21}. This is supplied dated curation derived from the OWR record, not an independently verified literature assessment. The supplied background expressly requests independent verification of time-sensitive status and cited claims.
\end{itemize}
\end{document}
